\chapter{Preamble}

\section*{Foreword}

The hardware platform employed throughout this project evolved in response to practical engineering constraints encountered at each stage of the work.

Initial exploration of the Dynamic Function eXchange (DFX) workflow — including Xilinx-provided tutorials, basic partial reconfiguration demonstrations, and familiarisation with the DFX Decoupler and ICAP primitives — was conducted on the \textbf{Digilent Arty A7-35T} (Xilinx Artix-7, XC7A35T). This board served as the primary vehicle for understanding the toolflow mechanics, constraint authoring, and bitstream generation procedures for 7-series partial reconfiguration.

Once the DFX fundamentals were established, development shifted to the \textbf{Digilent Arty A7-100T} (Xilinx Artix-7, XC7A100T) for the initial CFU Playground implementations and experimental runs. The larger logic capacity of the XC7A100T provided the additional resources required to accommodate the full LiteX-generated VexRiscv SoC alongside multiple reconfigurable CFU variants and the associated DFX infrastructure — a combination that proved too resource-intensive for the XC7A35T fabric.

For the final implementation and validation phase, the project transitioned to the \textbf{Digilent PYNQ-Z2} (Xilinx Zynq-7020). This decision was motivated by two considerations: first, the Zynq-7020 offered substantially greater programmable logic resources together with an integrated ARM Processing System, relaxing the area constraints that had limited design complexity on the Artix-7 platforms; second, the PYNQ framework's Python-based overlay and driver infrastructure enabled far more effective runtime monitoring, instrumentation, and interactive control of the reconfiguration process than was achievable through the bare-metal firmware and JTAG-only interfaces available on the Arty boards.

\vspace{0.4em}\hrule\vspace{0.4em}

\section{Introduction}

The proliferation of data-intensive workloads — spanning Deep Neural Networks (DNNs), real-time signal processing, and autonomous edge systems — has imposed increasingly stringent demands on computational hardware. General-purpose processors (GPPs), while versatile, frequently fail to deliver the throughput and energy efficiency required by these workloads within the area and power budgets characteristic of edge deployments. Application-Specific Integrated Circuits (ASICs), conversely, offer superior performance-per-watt but are inherently inflexible: once fabricated, their logic cannot be altered to accommodate evolving algorithmic requirements, changing operational contexts, or post-deployment functional corrections.

Field-Programmable Gate Arrays (FPGAs) occupy a well-established middle ground in this design-space trade-off, coupling the high-performance acceleration of custom digital logic with post-fabrication reconfigurability. Unlike ASICs, FPGA-based designs can be re-synthesised, re-placed, and re-routed to realise entirely new hardware functions on an existing silicon substrate. This intrinsic reconfigurability has positioned FPGAs as a critical compute platform across diverse application domains, including cloud acceleration, network function virtualisation, robotics, and embedded inference at the edge.

A defining capability of modern FPGA architectures, supported from the early Xilinx XC6200 family through the contemporary AMD Virtex UltraScale+ and 7-series devices, is \textit{Dynamic Function eXchange} (DFX), historically referred to as \textit{Dynamic Partial Reconfiguration} (DPR). DFX permits the selective modification of a designated region of the programmable logic fabric — termed a \textit{Reconfigurable Partition} (RP) — while the remainder of the design, the \textit{static partition}, continues to operate without interruption. At the hardware level, this is achieved by writing a \textit{partial bitstream} into the device's configuration memory through an internal primitive known as the \textit{Internal Configuration Access Port} (ICAP). The partial bitstream, comprising a synchronisation word, configuration frame addresses, and the corresponding frame data, reprograms only those configuration frames associated with the target RP, leaving all other frames — and hence all static logic — undisturbed.

The practical significance of DFX is considerable. In resource-constrained environments, it enables hardware time-multiplexing: a set of functionally distinct hardware accelerators can share a single physical region of the FPGA fabric, with only the currently required accelerator being instantiated at any given time. This approach reduces the total area footprint relative to a monolithic design in which all accelerators are implemented simultaneously. Furthermore, DFX decouples the compilation of individual reconfigurable modules from the compilation of the static infrastructure, thereby creating opportunities for parallel and incremental synthesis workflows that can substantially reduce overall build times — a significant consideration given that monolithic FPGA compilation routinely requires hours to complete.

The interest in DFX has intensified across multiple domains in recent years. In cloud and data-centre contexts, partially reconfigurable FPGA overlays enable multi-tenant virtualisation, wherein distinct users share a single physical device by occupying different reconfigurable regions. In edge and embedded systems, DFX facilitates adaptive compute: a single low-cost FPGA can dynamically reconfigure its accelerator logic in response to changing sensor inputs, mission phases, or operating-environment parameters. In the domain of hardware security, DFX has been explored both as a defence mechanism — through runtime diversification of cryptographic cores — and as a subject of vulnerability analysis.

In parallel with these advances in FPGA reconfiguration technology, the RISC-V Instruction Set Architecture (ISA) has emerged as a transformative paradigm for customisable processor design. As a free, open-standard ISA governed by a modular specification, RISC-V permits implementers to extend the base integer instruction set (RV32I) with application-specific custom instructions without incurring proprietary licensing costs. When realised as a soft-core processor on an FPGA — such as the VexRiscv core, an RV32I implementation written in SpinalHDL and deployed within the LiteX System-on-Chip (SoC) framework — RISC-V provides a flexible software execution environment directly on the reconfigurable fabric.

A primary mechanism for accelerating domain-specific workloads on such a soft-core processor is through the integration of \textit{Custom Function Units} (CFUs). The CFU Playground, an open-source framework developed by Google, facilitates the design, integration, and benchmarking of CFUs that attach to the VexRiscv core via a well-defined handshake interface comprising \texttt{cmd\_valid}, \texttt{cmd\_ready}, \texttt{rsp\_valid}, and \texttt{rsp\_ready} control signals, a 10-bit \texttt{function\_id} selector, and dual 32-bit operand inputs alongside a 32-bit result output. Each custom instruction encoded in the RISC-V CUSTOM0 opcode space is dispatched by the processor pipeline to the CFU, which executes the operation in hardware and returns the result within a deterministic number of clock cycles.

The convergence of these two paradigms — DFX-based partial reconfiguration and RISC-V custom instruction acceleration — presents a compelling architectural opportunity. Rather than statically binding a single CFU implementation to the processor at synthesis time, the CFU can be designated as a Reconfigurable Partition, allowing its hardware logic to be swapped at runtime via DFX. This approach enables a single physical RISC-V SoC to support a library of functionally diverse CFUs — for instance, transitioning from a byte-manipulation accelerator to a fixed-point multiply-accumulate unit — without re-synthesising or re-loading the entire bitstream. The processor, the memory subsystem, and the communication infrastructure remain operational throughout the reconfiguration process.

This project investigates the design, implementation, and validation of such a runtime-reconfigurable RISC-V co-processor architecture. The system is built upon the CFU Playground framework, targeting the Digilent Arty A7-35T development board (Xilinx Artix-7, XC7A35T) as the primary implementation platform, with baseline toolflow validation performed on the Arty A7-100T (XC7A100T) and supplementary cross-platform validation on the PYNQ-Z2 (Xilinx Zynq-7020). The AMD DFX IP suite, specifically the DFX Decoupler, is employed to isolate the reconfigurable CFU partition during bitstream loading, while the ICAP primitive and the STARTUPE2 End-of-Startup (EOS) signal are utilised for reconfiguration control and latency measurement, respectively.

\vspace{0.4em}\hrule\vspace{0.4em}

\section{Problem Statement and Objectives}

Despite its demonstrated potential, the practical adoption of Dynamic Partial Reconfiguration in FPGA-based systems remains constrained by several interrelated challenges.

\textbf{Toolflow complexity and accessibility.} State-of-the-art DPR toolflows — including Hierarchical Partial Reconfiguration (HPR), ZyPR, and RapidStream — have each advanced the field along specific dimensions, yet each presents significant barriers to practical adoption. HPR introduces a hierarchical page-based compilation strategy with recombinable Single, Double, and Quad PR pages interconnected by a Hoplite-BF Network-on-Chip (NoC), enabling parallel per-operator synthesis and adaptive resource allocation; however, reproduction of its toolflow was found to be hindered by strict Vivado version dependencies, Linux kernel compatibility requirements, and compute-resource constraints during place-and-route. RapidStream demonstrated order-of-magnitude reductions in compile time through latency-insensitive partitioning and parallel island-based compilation using RapidWright as the underlying infrastructure; however, the transition of RapidStream to a commercial product has resulted in the revocation of all public API access, rendering independent evaluation infeasible. ZyPR provides an abstracted, end-to-end build and runtime management flow for Zynq-based SoCs, incorporating the high-throughput ZyCAP PR controller with AXI-Lite control and AXI4 DMA-based bitstream transfer to the ICAP at up to 400 MB/s; however, dependent tooling requires an educational licence whose acquisition was unsuccessful, precluding hands-on evaluation. These experiences underscore the challenge of reproducing and extending existing DPR research due to proprietary, version-sensitive, and resource-intensive toolchain dependencies.

\textbf{Hardware controller requirements for non-Zynq platforms.} The majority of existing DPR controller architectures — including ZyCAP and the standard Linux FPGA Manager interface (which relies on the Processor Configuration Access Port, PCAP, on Zynq devices) — assume the availability of an integrated ARM Processing System (PS) for orchestrating bitstream transfers. On pure-FPGA platforms such as the Artix-7 family, no such PS exists. Reconfiguration must therefore be managed entirely from within the programmable logic (PL) fabric, using only the ICAP primitive and locally available memory resources. The design of a PL-autonomous reconfiguration controller that can fetch, validate, and write partial bitstreams to the ICAP without host intervention remains an open engineering challenge for 7-series devices.

\textbf{Integration with soft-core processor ecosystems.} While RISC-V-based DPR approaches have been explored in the literature — notably RV-CAP, which demonstrated ICAP-based reconfiguration orchestrated by a VexRiscv soft core, and AMPER-X, which employed DFX to dynamically alter the precision of arithmetic units attached to a RISC-V core — the specific integration of DFX with the CFU Playground framework has not been previously reported. The CFU Playground's standardised CFU interface and its tightly coupled relationship with the VexRiscv pipeline present both an opportunity for clean architectural partitioning and a set of design constraints around handshake signal isolation, clock domain management, and decoupler placement that must be addressed for correct DFX operation.

In light of these challenges, the objectives of the present work are defined as follows:

\begin{enumerate}
  \item \textbf{Literature survey of state-of-the-art DPR toolflows and architectures.} A comprehensive review of existing approaches to DPR — encompassing HPR, ZyPR/ZyCAP, RapidStream, DORA, RV-CAP, AMPER-X, DML, and ACNNE — is conducted to identify the architectural principles, compilation strategies, and controller designs that inform the proposed system. Where hands-on evaluation was not feasible due to licensing, version, or resource constraints, the analysis is restricted to a literature-based assessment of the published results.
  \item \textbf{Understanding of hardware controller requirements.} The ICAP primitive, the DFX Decoupler IP, and the STARTUPE2 configuration primitive on 7-series FPGAs are studied in detail to establish the hardware infrastructure necessary for controlled, non-disruptive partial reconfiguration on the Arty A7-35T.
  \item \textbf{Design and implementation of a runtime-reconfigurable RISC-V co-processor.} A system architecture is proposed and implemented in which the \texttt{cfu\_compute} module — the functional core of the CFU — is designated as a Reconfigurable Partition within a DFX-enabled Vivado project. The static partition comprises the VexRiscv SoC (generated by LiteX and wrapped in the \texttt{digilent\_arty} module), the DFX Decoupler (instantiated within \texttt{cfu.v}), the ICAP controller infrastructure (within \texttt{system\_wrapper}), and the reconfiguration latency measurement logic (\texttt{recon\_counter.v} utilising the STARTUPE2 EOS signal). Multiple CFU implementations — including \texttt{example.v} (byte sum, byte swap, and bit reverse operations), \texttt{donut.v} (signed fixed-point multiply and multiply-shift-right), and \texttt{dse\_template.v} — are synthesised as distinct Reconfigurable Modules (RMs), each producing an independent partial bitstream.
  \item \textbf{Validation across multiple hardware platforms and reconfiguration methods.} The reconfiguration flow is validated using JTAG-based bitstream delivery on both the Arty A7-35T and the PYNQ-Z2, and additionally via PCAP-based reconfiguration through the Zynq PS on the PYNQ-Z2. Reconfiguration latency is measured using an Integrated Logic Analyser (ILA) core connected to the \texttt{recon\_counter} output, which increments while the STARTUPE2 EOS signal is de-asserted during the reconfiguration window.
  \item \textbf{Identification of a roadmap toward fully autonomous reconfiguration.} The current implementation employs host-driven, JTAG-based partial bitstream loading. The intended progression toward autonomous reconfiguration — first via SPI/QSPI flash fetch using the DFX Controller IP, and subsequently via DDR-backed runtime storage with DMA-based bitstream transfer to the ICAP — is documented as a defined engineering roadmap for future work.
\end{enumerate}